\documentclass[11pt,a4paper]{article}

\usepackage[T1]{fontenc}
\usepackage[utf8]{inputenc}
\usepackage[ngerman]{babel}

\usepackage{amsmath,amssymb,amsthm}
\usepackage[a4paper,margin=2.5cm]{geometry}

\title{Chebyshev-artige Restklassenasymmetrie als mögliche Quelle einer Schalenbalance}
\author{Kollaborative Ausarbeitung}
\date{\today}

\newtheorem{satz}{Satz}
\newtheorem{lemma}[satz]{Lemma}
\newtheorem{korollar}[satz]{Korollar}
\theoremstyle{definition}
\newtheorem{definition}[satz]{Definition}
\newtheorem{hypothese}[satz]{Hypothese}
\newtheorem{beispiel}[satz]{Beispiel}
\theoremstyle{remark}
\newtheorem{bemerkung}[satz]{Bemerkung}

\begin{document}
	\maketitle
	
	\begin{abstract}
		Wir formulieren ein Modell zur Untersuchung einer numerisch beobachteten Balance zwischen
		einem glatten Anteil und einem energetischen Anteil natürlicher Zahlen. Dazu werden zu einer
		wachsenden Primschale \(P_k\) der \(P_k\)-glatte Anteil \(S_k(n)\), der schalenreduzierte Rest
		\(R_k(n)\) und dessen energetischer Anteil \(E_k(n)\) in der Restklasse \(1 \pmod{12}\)
		definiert. Aus diesen Größen werden logarithmische und symmetrisch normierte Balancefunktionen
		abgeleitet. Der Artikel formuliert die Arbeitshypothese, dass die beobachtete Schalenbalance
		nicht allein ein Effekt der Glättung ist, sondern zumindest teilweise aus einer
		Chebyshev-artigen Restklassenasymmetrie modulo \(12\) stammt, die durch die multiplikative
		Struktur der Primfaktorzerlegung verstärkt sichtbar wird. Abschließend wird ein präzises
		numerisches Testprogramm formuliert.
	\end{abstract}
	
	\section{Einleitung}
	
	Der klassische Chebyshev-Bias beschreibt die Beobachtung, dass Primzahlen auf endlichen Skalen
	nicht immer gleichmäßig auf Restklassen verteilt erscheinen, obwohl asymptotisch Gleichverteilung
	erwartet wird. Typischerweise betrifft dies Zählfunktionen der Form
	\[
	\pi(x;q,a),
	\]
	also die Anzahl der Primzahlen \(\le x\), die in einer festen Restklasse \(a \pmod q\) liegen.
	
	Im vorliegenden Zusammenhang geht es nicht um solche additiven Zählgrößen, sondern um ein
	multiplikatives Modell. Für eine wachsende Primschale \(P_k\) wird zunächst der glatte Anteil
	einer Zahl abgetrennt. Im verbleibenden Rest wird dann der Anteil der Primfaktorpotenzen in der
	Klasse \(1 \pmod{12}\) isoliert. Die numerische Beobachtung ist, dass aus diesen Größen gebildete
	Balancefunktionen auf bestimmten Grundmengen eine auffällige Drift in Richtung einer neutralen
	Lage zeigen.
	
	Ziel dieses Artikels ist es, diese Beobachtung präzise zu formulieren und die Hypothese
	auszusprechen, dass hier nicht nur ein Glättungseffekt, sondern eine verfeinerte,
	multiplikative Erscheinungsform einer Chebyshev-artigen Restklassenasymmetrie vorliegt.
	
	\section{Schalen, Reste und energetische Anteile}
	
	\begin{definition}[Primschalen]
		Für \(k\in\mathbb{N}\) sei
		\[
		P_k=\{2,3,5,\dots,p_k\}
		\]
		die Menge der ersten \(k\) Primzahlen.
	\end{definition}
	
	\begin{definition}[Glatter Anteil und Rest]
		Für \(n\in\mathbb{N}\) definieren wir den \(P_k\)-glatten Anteil durch
		\[
		S_k(n):=\prod_{p\in P_k} p^{v_p(n)}
		\]
		und den zugehörigen Rest durch
		\[
		R_k(n):=\frac{n}{S_k(n)}.
		\]
	\end{definition}
	
	\begin{bemerkung}
		Der Anteil \(S_k(n)\) enthält genau die Primfaktorpotenzen von \(n\), deren Primzahlen in der
		Schale \(P_k\) liegen. Der Rest \(R_k(n)\) ist zu allen Primzahlen aus \(P_k\) teilerfremd.
	\end{bemerkung}
	
	\begin{definition}[Energetischer Anteil]
		Der energetische Anteil des Restes \(R_k(n)\) sei definiert durch
		\[
		E_k(n):=\prod_{\substack{q^a\parallel R_k(n)\\ q\equiv 1\;(\mathrm{mod}\;12)}} q^a.
		\]
	\end{definition}
	
	\begin{definition}[Restklassenanteile des Restes]
		Zusätzlich definieren wir die drei weiteren Restklassenanteile
		\[
		A_k(n):=\prod_{\substack{q^a\parallel R_k(n)\\ q\equiv 5\;(\mathrm{mod}\;12)}} q^a,
		\]
		\[
		B_k(n):=\prod_{\substack{q^a\parallel R_k(n)\\ q\equiv 7\;(\mathrm{mod}\;12)}} q^a,
		\]
		\[
		C_k(n):=\prod_{\substack{q^a\parallel R_k(n)\\ q\equiv 11\;(\mathrm{mod}\;12)}} q^a.
		\]
	\end{definition}
	
	\begin{bemerkung}
		Formal zerfällt der Rest \(R_k(n)\) in die vier modulo-\(12\)-Komponenten
		\[
		R_k(n)=E_k(n)\,A_k(n)\,B_k(n)\,C_k(n).
		\]
		Da \(R_k(n)\) bereits zu allen Primzahlen der Schale \(P_k\) teilerfremd ist, treten darin keine
		weiteren Primfaktoren aus \(P_k\) auf.
	\end{bemerkung}
	
	\section{Balancegrößen}
	
	\begin{definition}[Logarithmische Balance]
		Wir definieren
		\[
		L_k(n):=\log E_k(n)-\log S_k(n).
		\]
	\end{definition}
	
	\begin{definition}[Symmetrisch normierte Balance]
		Die zugehörige normierte Balancefunktion sei
		\[
		G_k(n):=\frac12+\frac1\pi\arctan\!\bigl(L_k(n)\bigr).
		\]
	\end{definition}
	
	\begin{bemerkung}
		Die Funktion \(G_k(n)\) nimmt Werte im offenen Intervall \((0,1)\) an.
		Dabei bedeutet
		\[
		G_k(n)=\frac12
		\]
		eine exakte Balance zwischen energetischem Anteil und glattem Anteil. Ferner gilt:
		\[
		G_k(n)>\frac12
		\quad\Longleftrightarrow\quad
		E_k(n)>S_k(n),
		\]
		\[
		G_k(n)<\frac12
		\quad\Longleftrightarrow\quad
		E_k(n)<S_k(n).
		\]
	\end{bemerkung}
	
	\begin{definition}[Mittelwerte]
		Für \(N\in\mathbb{N}\) definieren wir die Mittelwerte
		\[
		M_k(N):=\frac{1}{N}\sum_{n\le N}L_k(n),
		\qquad
		H_k(N):=\frac{1}{N}\sum_{n\le N}G_k(n).
		\]
		Auf der restringierten Grundmenge
		\[
		\mathcal N_6(N):=\{n\le N:\gcd(n,6)=1\}
		\]
		definieren wir analog
		\[
		M_k^\ast(N):=
		\frac{1}{|\mathcal N_6(N)|}
		\sum_{\substack{n\le N\\ \gcd(n,6)=1}}L_k(n),
		\]
		\[
		H_k^\ast(N):=
		\frac{1}{|\mathcal N_6(N)|}
		\sum_{\substack{n\le N\\ \gcd(n,6)=1}}G_k(n).
		\]
	\end{definition}
	
	\section{Chebyshev-artige Restklassenasymmetrie}
	
	\begin{definition}[Chebyshev-artige Restklassenasymmetrie modulo \(12\)]
		Wir sagen, dass auf der Ebene der Primfaktorbeiträge eine Chebyshev-artige
		Restklassenasymmetrie modulo \(12\) vorliegt, wenn die Restklassen
		\[
		1,\ 5,\ 7,\ 11 \pmod{12}
		\]
		auf endlichen Skalen nicht symmetrisch zu den logarithmischen Beiträgen der Primfaktoren des
		Restes \(R_k(n)\) beitragen, das heißt wenn die Größen
		\[
		\log E_k(n),\quad \log A_k(n),\quad \log B_k(n),\quad \log C_k(n)
		\]
		im Mittel systematisch verschieden verteilt sind.
	\end{definition}
	
	\begin{bemerkung}
		Diese Definition betrifft nicht die klassische Primzahlzählung, sondern ein
		multiplikativ-logarithmisches Modell. Sie ist also nicht identisch mit dem klassischen
		Chebyshev-Bias, sondern als dessen verfeinerte, auf Primfaktorprodukte bezogene Variante zu
		verstehen.
	\end{bemerkung}
	
	\section{Heuristische Verstärkung durch Multiplikativität}
	
	\begin{lemma}[Heuristische Verstärkung]
		Liegt auf den Resten \(R_k(n)\) eine Chebyshev-artige Restklassenasymmetrie modulo \(12\) vor,
		so ist heuristisch zu erwarten, dass die Mittelwerte
		\[
		M_k(N)=\frac{1}{N}\sum_{n\le N}L_k(n)
		\]
		und damit auch
		\[
		H_k(N)=\frac{1}{N}\sum_{n\le N}G_k(n)
		\]
		systematisch von der neutralen Lage \(0\) beziehungsweise \(\frac12\) abweichen.
	\end{lemma}
	
	\begin{proof}
		Die Begründung ist heuristisch. Restklassenasymmetrien werden hier nicht additiv gezählt,
		sondern über Produkte von Primfaktorpotenzen zu logarithmischen Summen verarbeitet. Dadurch
		können kleine Unterschiede in Häufigkeit oder Größe von Primfaktoren aus der Klasse
		\(1 \pmod{12}\) gegenüber den Klassen \(5,7,11 \pmod{12}\) in den Größen \(L_k(n)\) und
		\(G_k(n)\) verstärkt sichtbar werden.
	\end{proof}
	
	\section{Arbeitshypothese}
	
	\begin{hypothese}[Schalenbalance als Restklasseneffekt]
		Die numerisch beobachtete Schalenbalance der Größen \(G_k(n)\) ist plausibel als
		multiplikative, schalenmodulierte Erscheinungsform einer Chebyshev-artigen
		Restklassenasymmetrie modulo \(12\) zu deuten. Genauer:
		\begin{enumerate}
			\item Die Schale \(S_k(n)\) entfernt systematisch kleine Primfaktoren.
			\item Der Rest \(R_k(n)\) trägt die eigentliche restklassenspezifische Primfaktorstruktur.
			\item Der energetische Anteil \(E_k(n)\) selektiert innerhalb dieses Restes genau die Klasse
			\[
			1\pmod{12}.
			\]
			\item Jede Asymmetrie dieser Klasse relativ zu den Klassen
			\[
			5,\ 7,\ 11 \pmod{12}
			\]
			kann in \(L_k(n)\) und \(G_k(n)\) multiplikativ verstärkt erscheinen.
		\end{enumerate}
	\end{hypothese}
	
	\begin{satz}[Heuristische Konsequenz]
		Unter der vorstehenden Hypothese ist es plausibel, dass die Drift der Mittelwerte
		\[
		H_k(N)\quad\text{und}\quad H_k^\ast(N)
		\]
		mit wachsender Primschale nicht bloß ein reiner Glättungseffekt ist, sondern zumindest teilweise
		durch eine tieferliegende Restklassenasymmetrie verursacht wird.
	\end{satz}
	
	\begin{proof}
		Dies ist keine strenge Deduktion, sondern eine strukturierte Schlussfolgerung aus der Hypothese
		und dem heuristischen Lemma. Der Punkt ist, dass \(E_k(n)\) eine ausgezeichnete Restklasse
		selektiert, während \(S_k(n)\) nur die kleinen Schalenfaktoren abzieht. Dadurch können sich
		systematische Unterschiede der Restklassen im schalenreduzierten Rest in den Balancegrößen
		niederschlagen.
	\end{proof}
	
	\section{Numerisches Testprogramm}
	
	\begin{korollar}[Prüfbare Version der Hypothese]
		Zur numerischen Prüfung der Hypothese genügt es, für wachsendes \(N\) und feste Schalen \(P_k\)
		die Mittelwerte
		\[
		\frac{1}{N}\sum_{n\le N}\log E_k(n),\qquad
		\frac{1}{N}\sum_{n\le N}\log A_k(n),\qquad
		\frac{1}{N}\sum_{n\le N}\log B_k(n),\qquad
		\frac{1}{N}\sum_{n\le N}\log C_k(n)
		\]
		sowie die daraus abgeleiteten Balancegrößen \(M_k(N)\), \(H_k(N)\), \(M_k^\ast(N)\),
		\(H_k^\ast(N)\) miteinander zu vergleichen.
	\end{korollar}
	
	\begin{proof}
		Wenn die beobachtete Schalenbalance tatsächlich von einer restklassenspezifischen Asymmetrie
		getrieben wird, dann sollte sich dies bereits auf der Ebene der vier logarithmischen
		Restklassenanteile zeigen. Bleibt der Anteil \(E_k\) systematisch gegen \(A_k,B_k,C_k\)
		verschoben, so spricht dies für einen restklassentheoretischen Ursprung des Effekts.
	\end{proof}
	
	\begin{bemerkung}
		Insbesondere ist zu unterscheiden zwischen
		\begin{enumerate}
			\item der Vollmenge aller Zahlen \(n\le N\),
			\item der restringierten Menge \(\{n\le N:\gcd(n,6)=1\}\).
		\end{enumerate}
		Wenn die Schalenbalance auf der zweiten Menge systematisch näher an eine neutrale Lage
		\(\frac12\) heranrückt, auf der ersten aber nicht, dann spricht das für einen strukturellen
		Zusammenhang zwischen der modulo-\(12\)-Zerlegung und der Entfernung kleiner glatter Faktoren.
	\end{bemerkung}
	
	\section{Diskussion}
	
	Der klassische Chebyshev-Bias ist eine Aussage über Zählfunktionen von Primzahlen in
	Restklassen. Das hier betrachtete Modell arbeitet hingegen mit Primfaktorprodukten,
	logarithmischen Summen und glatten Schalen. Darum ist die passende Aussage nicht
	\[
	\text{Schalenbalance}=\text{Chebyshev-Bias},
	\]
	sondern vielmehr
	\[
	\text{Schalenbalance}
	=
	\text{Chebyshev-artige Restklassenasymmetrie}
	+
	\text{multiplikative Schalengeometrie}.
	\]
	
	Gerade diese Kombination macht das Modell interessant. Die Glättung durch \(S_k(n)\) entfernt
	kleine Primfaktoren, der Rest \(R_k(n)\) trägt die nichttriviale Restklassenstruktur, und die
	Selektion auf \(E_k(n)\) hebt eine spezielle Restklasse hervor. Auf diese Weise könnte aus einem
	schwachen restklassentheoretischen Bias ein deutlich sichtbarer Effekt in den Größen
	\(L_k(n)\) und \(G_k(n)\) entstehen.
	
	\section{Schluss}
	
	Wir haben ein präzises Rahmenmodell formuliert, in dem die numerisch beobachtete Balance zwischen
	glattem Anteil und energetischem Anteil mathematisch erfasst werden kann. Die zentrale
	Arbeitshypothese lautet, dass diese Balance zumindest teilweise durch eine Chebyshev-artige
	Restklassenasymmetrie modulo \(12\) erklärt werden kann, die durch die multiplikative Struktur
	der Primfaktorzerlegung verstärkt wird.
	
	Der Artikel liefert keinen Beweis dieser Hypothese, wohl aber
	\begin{enumerate}
		\item eine saubere formale Definition der beteiligten Größen,
		\item ein heuristisches Verstärkungslemma,
		\item eine präzise formulierte Hypothese,
		\item und ein konkretes numerisches Testprogramm.
	\end{enumerate}
	
	Damit liegt eine belastbare Grundlage vor, um die beobachtete Schalenbalance weiter theoretisch
	und numerisch zu untersuchen.
	
\end{document}